\documentclass[a4paper,12pt]{article}
\usepackage[T2A]{fontenc}
\usepackage[utf8]{inputenc}
\usepackage[russian]{babel}
\usepackage{amsmath,amssymb,mathtools}
\renewcommand{\familydefault}{\sfdefault}
\usepackage{icomma}
\usepackage{geometry}
\geometry{left=18mm,right=18mm,top=18mm,bottom=19mm}
\usepackage{microtype}
\usepackage{tikz}
\usetikzlibrary{calc,arrows.meta,positioning,shapes.geometric}
\usepackage{tabularx,booktabs,longtable,array}
\usepackage{enumitem}
\usepackage{hyperref}
\usepackage{parskip}
\usepackage{fancyhdr}
\usepackage{setspace}
\usepackage{xcolor}

\definecolor{accent}{HTML}{2B6F73}
\definecolor{darkaccent}{HTML}{19484B}
\definecolor{textgray}{HTML}{2E3335}
\definecolor{softgray}{HTML}{6D777A}
\definecolor{linegray}{HTML}{D8DEDF}

\hypersetup{colorlinks=true,linkcolor=darkaccent,urlcolor=darkaccent}
\onehalfspacing
\setlength{\parindent}{0pt}
\setlist[itemize]{leftmargin=6mm,itemsep=2pt,topsep=3pt}
\setlist[enumerate]{leftmargin=8mm,itemsep=4pt,topsep=4pt}
\setlength{\headheight}{14.5pt}
\pagestyle{fancy}
\fancyhf{}
\fancyhead[L]{\small\color{softgray}София Хоменко}
\fancyhead[R]{\small\color{softgray}06.10.26}
\fancyfoot[C]{\small\color{softgray}\thepage}
\renewcommand{\headrulewidth}{0pt}
\renewcommand{\footrulewidth}{0pt}

\newcommand{\topic}[1]{%
  \vspace{0.45em}\noindent{\Large\bfseries\color{darkaccent}#1}\par
  \vspace{0.25em}\noindent\color{linegray}\rule{\linewidth}{0.6pt}\color{textgray}\vspace{0.45em}
}
\newcommand{\key}[1]{\textbf{\color{darkaccent}#1}}
\newcommand{\formula}[1]{\[\boxed{\displaystyle #1}\]}
\newcommand{\warn}{\textbf{\color{darkaccent}Важно: }}

\begin{document}
\color{textgray}

% ---------- TITLE ----------
\thispagestyle{empty}
\begin{center}
\vspace*{24mm}
{\Large\color{softgray}Математика • ОГЭ}\par
\vspace{10mm}
{\Huge\bfseries\color{darkaccent}Чек-лист}\par
\vspace{5mm}
{\LARGE\bfseries Задачи про террасы:\par Пифагор, уклон, площади и проценты}\par
\vspace{15mm}
{\large София Хоменко}\par
\vspace{4mm}
{\normalsize 6 октября 2026 года}\par
\vfill
{\small Лёвин Артём Александрович эксклюзивно для Софии Хоменко}\par
\vspace{12mm}
\begin{tikzpicture}[x=1cm,y=1cm]
  \draw[accent,line width=1.2pt]
  (0,0) .. controls (3.2,0.8) and (7.3,-0.8) .. (10.5,0)
        .. controls (13.1,0.6) and (15.6,0.2) .. (17.3,-0.15);
\end{tikzpicture}
\end{center}
\clearpage

% ---------- MAIN CHECKLIST ----------
\topic{1. Сначала переведите текст задачи на язык величин}

В задачах про склон и террасы удобно сразу ввести четыре обозначения:
\begin{itemize}
  \item \(h\) --- перепад высот, вертикальный катет;
  \item \(a\) --- горизонтальная проекция склона, второй катет;
  \item \(l\) --- длина склона, гипотенуза;
  \item \(w\) --- ширина участка.
\end{itemize}

Эти величины образуют прямоугольный треугольник, поэтому основная геометрическая связь:
\formula{l^2=h^2+a^2,\qquad l=\sqrt{h^2+a^2}.}

\textbf{Пример из занятия.} При \(h=15\) м и \(a=112\) м:
\[
l^2=15^2+112^2=225+12544=12769,
\]
\[
l=\sqrt{12769}=113\text{ м}.
\]

\warn перед вычислениями подпишите, какая длина является вертикальной, какая горизонтальной, а какая --- длиной склона. Это снимает большую часть ошибок в последующих пунктах.

\topic{2. Как быстро распознать квадрат большого числа}

Если таблица квадратов заканчивается слишком рано, используйте два признака.

\begin{enumerate}
  \item \key{Определите десяток корня.} Сравните число с квадратами круглых десятков.
  \item \key{Посмотрите на последнюю цифру квадрата.} Она резко сужает число кандидатов.
\end{enumerate}

Полезные пары для последней цифры корня:
\[
\begin{array}{c|cccccc}
\text{последняя цифра квадрата} & 0 & 1 & 4 & 5 & 6 & 9\\
\hline
\text{последняя цифра корня} & 0 & 1\text{ или }9 & 2\text{ или }8 & 5 & 4\text{ или }6 & 3\text{ или }7
\end{array}
\]

\textbf{Пример.} Найдём \(\sqrt{5184}\).
\[
70^2=4900<5184<6400=80^2.
\]
Значит, корень лежит между 70 и 80. Число оканчивается на 4, поэтому единицы корня --- 2 или 8. Кандидаты: \(72\) и \(78\). Проверяем:
\[
72^2=5184.
\]
Следовательно, \(\sqrt{5184}=72\).

\warn последняя цифра даёт только кандидатов. Финальная проверка возведением в квадрат обязательна.

\topic{3. Уклон в процентах: это не градусы}

В прикладной задаче «угол склона в процентах» может быть задан через тангенс:
\formula{p=\tg\alpha\cdot100\%=\frac{h}{a}\cdot100\%.}

Для \(h=15\) м и \(a=112\) м:
\[
p=\frac{15}{112}\cdot100\%\approx13{,}392857\%\approx13{,}4\%.
\]

Если требуется проверить ограничение, например «уклон не более \(50\%\)», сравнивайте уже полученное процентное значение с \(50\%\).

\warn число \(50\%\) в таком условии не означает угол \(50^\circ\). Следуйте определению, данному в тексте задачи.

\clearpage

\topic{4. Площадь участка до устройства террас}

До террасирования посевная поверхность лежит вдоль склона. Если участок прямоугольный, его площадь равна
\formula{S_{\text{до}}=w\cdot l.}

Для данных занятия \(w=30\) м и \(l=113\) м:
\[
S_{\text{до}}=30\cdot113=3390\text{ м}^2.
\]

\topic{5. Площадь после устройства равных террас}

Пусть склон разбит на \(n\) одинаковых горизонтальных ступеней. Горизонтальная длина одной ступени равна
\[
\frac{a}{n}.
\]
Площадь одной ступени:
\[
S_1=w\cdot\frac{a}{n}.
\]
Всего ступеней \(n\), поэтому
\[
S_{\text{после}}=n\cdot S_1
=n\cdot w\cdot\frac{a}{n}
=w\cdot a.
\]

\formula{S_{\text{после}}=w\cdot a.}

Для \(w=30\) м и \(a=112\) м:
\[
S_{\text{после}}=30\cdot112=3360\text{ м}^2.
\]

\textbf{Главный вывод:} при равных террасах число ступеней в итоговой формуле сокращается. Оно нужно для понимания длины одной ступени, но суммарная горизонтальная площадь определяется произведением \(w\cdot a\).

\warn не делите площадь всего участка на число ступеней. Делится горизонтальная длина \(a\), после чего площадь одной ступени умножается на количество ступеней.

\topic{6. На сколько процентов уменьшилась площадь}

Если величина уменьшилась с \(A_{\text{до}}\) до \(A_{\text{после}}\), процент уменьшения считается от исходного значения:
\formula{q=\frac{A_{\text{до}}-A_{\text{после}}}{A_{\text{до}}}\cdot100\%.}

Для площадей занятия:
\[
q=\frac{3390-3360}{3390}\cdot100\%
=\frac{30}{3390}\cdot100\%
\approx0{,}885\%\approx0{,}9\%.
\]

\warn в знаменателе стоит то, что было \emph{до изменения}. При вопросе «на сколько процентов сократилось» удобно сразу вычитать меньшее значение из большего и получать положительный процент уменьшения.

\topic{7. Урожай и потеря массы при обработке}

Если с каждого квадратного метра получают \(r\) граммов продукта, то масса урожая до обработки:
\formula{M_{\text{сыр}}=S\cdot r.}

Если при обработке теряется \(k\%\) массы, остаётся
\[
100\%-k\%.
\]
В виде десятичного множителя:
\[
1-\frac{k}{100}.
\]
Поэтому масса готового продукта в килограммах:
\formula{M_{\text{готов}}=\frac{S\cdot r\left(1-\frac{k}{100}\right)}{1000}.}

\textbf{Пример из занятия.} При \(S=3360\text{ м}^2\), урожайности \(600\text{ г/м}^2\) и потере \(15\%\):
\[
M_{\text{сыр}}=3360\cdot600=2\,016\,000\text{ г},
\]
\[
M_{\text{готов}}=2\,016\,000\cdot0{,}85=1\,713\,600\text{ г}=1713{,}6\text{ кг}.
\]

\warn потеря \(15\%\) означает, что остаётся \(85\%\), то есть нужно умножить на \(0{,}85\).

% ---------- FLOWCHART ----------
\clearpage
\thispagestyle{plain}
\begin{center}
{\Large\bfseries\color{darkaccent}Алгоритм: задача про склон и террасы}\par
\vspace{8mm}
\end{center}

\tikzset{
  startstop/.style={ellipse,draw=accent,line width=0.9pt,align=center,minimum width=45mm,minimum height=10mm},
  process/.style={rounded corners=3pt,draw=accent,line width=0.9pt,align=center,text width=120mm,minimum height=11mm,inner sep=3.5mm},
  arrow/.style={-{Latex[length=3mm]},line width=0.9pt,color=darkaccent}
}
\begin{center}
\begin{tikzpicture}[node distance=8mm]
  \node[startstop] (s) {Прочитать условие и выписать данные};
  \node[process,below=of s] (vars) {Обозначить \(h\), \(a\), \(l\), \(w\). Отдельно выписать число террас, урожайность и проценты, если они даны};
  \node[process,below=of vars] (pyth) {Если нужна длина склона: \(l=\sqrt{h^2+a^2}\)};
  \node[process,below=of pyth] (slope) {Если нужен уклон: \(p=\dfrac{h}{a}\cdot100\%\)};
  \node[process,below=of slope] (areas) {Площади: \(S_{\text{до}}=wl\), \(S_{\text{после}}=wa\)};
  \node[process,below=of areas] (percent) {Если спрашивают уменьшение: \(q=\dfrac{S_{\text{до}}-S_{\text{после}}}{S_{\text{до}}}\cdot100\%\)};
  \node[process,below=of percent] (yield) {Если нужен урожай: площадь \(\to\) граммы \(\to\) учёт потери массы \(\to\) перевод в килограммы};
  \node[startstop,below=of yield] (f) {Округлить только до точности из условия и записать ответ};

  \draw[arrow] (s) -- (vars);
  \draw[arrow] (vars) -- (pyth);
  \draw[arrow] (pyth) -- (slope);
  \draw[arrow] (slope) -- (areas);
  \draw[arrow] (areas) -- (percent);
  \draw[arrow] (percent) -- (yield);
  \draw[arrow] (yield) -- (f);
\end{tikzpicture}
\end{center}
\clearpage

% ---------- TRAINING ----------
\topic{Тренировочный блок — Путь А}

\textbf{10 заданий.} Идите по нарастающей сложности. В задачах на проценты округляйте ответ до десятых, если отдельно не сказано иное.

\begin{enumerate}
  \item Не пользуясь таблицей квадратов, найдите \(\sqrt{5184}\). В решении укажите, между квадратами каких круглых десятков лежит число и какие две последние цифры корня возможны.

  \item Участок имеет ширину \(24\) м. Перепад высот равен \(20\) м, горизонтальная проекция склона --- \(99\) м. Найдите длину склона и площадь участка до устройства террас.

  \item Перепад высот равен \(18\) м, горизонтальная проекция склона --- \(80\) м. Найдите уклон в процентах.

  \item Террасы разрешено строить при уклоне не более \(20\%\). Для участка перепад высот равен \(14\) м, горизонтальная проекция --- \(75\) м. Определите уклон и решите, выполняется ли ограничение.

  \item Участок шириной \(24\) м имеет горизонтальную проекцию склона \(99\) м. Склон разделили на \(9\) равных террас. Найдите горизонтальную длину одной ступени, площадь одной ступени и суммарную посевную площадь после устройства террас.
\end{enumerate}

\clearpage
\topic{Тренировочный блок — продолжение}
\begin{enumerate}
  \setcounter{enumi}{5}
  \item Для участка шириной \(24\) м длина склона равна \(101\) м, а его горизонтальная проекция --- \(99\) м. На сколько процентов уменьшилась посевная площадь после устройства террас? Ответ округлите до десятых.

  \item После устройства террас посевная площадь равна \(2376\text{ м}^2\). С одного квадратного метра получают \(550\) г бурого риса. При шлифовке теряется \(12\%\) массы. Сколько килограммов белого риса получится? Ответ округлите до десятых килограмма.

  \item Участок шириной \(28\) м имеет перепад высот \(12\) м и горизонтальную проекцию \(35\) м. Найдите длину склона, площадь до террасирования, площадь после террасирования и процент уменьшения площади. Процент округлите до десятых.

  \item Для участка перепад высот равен \(21\) м, горизонтальная проекция --- \(72\) м. а) Найдите длину склона. б) Найдите уклон в процентах. в) Разрешено ли устройство террас, если допустимый уклон не превышает \(30\%\)?

  \item Участок шириной \(28\) м имеет горизонтальную проекцию склона \(35\) м. После устройства террас с каждого квадратного метра собирают \(640\) г бурого риса. При шлифовке теряется \(15\%\) массы. Сколько килограммов белого риса получится со всего участка? Ответ округлите до десятых килограмма.
\end{enumerate}

\clearpage

% ---------- ANSWERS ----------
\topic{Ответы}

\renewcommand{\arraystretch}{1.35}
\begin{longtable}{>{\raggedright\arraybackslash}p{11mm} >{\raggedright\arraybackslash}p{0.80\linewidth}}
\toprule
\textbf{№} & \textbf{Ответ} \\
\midrule
\endhead
1 & \(72\). Число лежит между \(70^2\) и \(80^2\); возможные единицы корня --- \(2\) или \(8\). \\
2 & \(l=101\) м; \(S_{\text{до}}=2424\text{ м}^2\). \\
3 & \(22{,}5\%\). \\
4 & \(18{,}7\%\); ограничение выполняется. \\
5 & \(11\) м; \(264\text{ м}^2\); \(2376\text{ м}^2\). \\
6 & \(2{,}0\%\). \\
7 & \(1150{,}0\) кг. \\
8 & \(l=37\) м; \(S_{\text{до}}=1036\text{ м}^2\); \(S_{\text{после}}=980\text{ м}^2\); уменьшение \(5{,}4\%\). \\
9 & \(75\) м; \(29{,}2\%\); да, разрешено. \\
10 & \(533{,}1\) кг. \\
\bottomrule
\end{longtable}

\vfill
\begin{center}
{\small\color{softgray}Лёвин Артём Александрович эксклюзивно для Софии Хоменко}
\end{center}

\end{document}
